\section{Chain fission audit worked row}

The fission audit uses the same glycine--alanine row and checks that the declared split recovers the route unit.

\subsection{Declare the split}

For the split after the first component, the audit identity is
\[
\boxed{\mathbf n(\mathrm{Gly})+\mathbf n(\mathrm{Ala})
=
\mathbf n(\mathrm{Gly\mbox{-}Ala})+\mathbf n(H_2O).}
\]

\subsection{Compute both sides}

Left side:
\[
(2,5,1,2,0,0)+(3,7,1,2,0,0)=(5,12,2,4,0,0).
\]
Right side:
\[
(5,10,2,3,0,0)+(0,2,0,1,0,0)=(5,12,2,4,0,0).
\]
Therefore
\[
\Delta_{\mathrm{CHNOPS}}=(0,0,0,0,0,0),
\qquad
\sum |\Delta_{\mathrm{CHNOPS}}|=0.
\]

\subsection{Audit row}

The audit row is recorded under \texttt{chain\_GA\_peptide\_seed} in
\begin{center}
\small\path{manual-2/data/molecular/chain_fission_audit.csv}
\end{center}
The audit status is \texttt{passed\_exact\_route\_recovery}.
